\documentclass[11pt,a4paper]{article}

%% Motor: xelatex (o sube el .tex a overleaf.com)
\usepackage{fontspec}
\usepackage{polyglossia}
\setmainlanguage{spanish}
\usepackage{amsmath,amssymb,amsthm}
\usepackage{geometry}
\usepackage{multicol}
\usepackage{enumitem}
\usepackage{fancyhdr}
\usepackage{lastpage}
\usepackage{xcolor}
\usepackage[most]{tcolorbox}
\usepackage{hyperref}
\usepackage{booktabs}

\geometry{top=2cm,bottom=2.2cm,left=2cm,right=2cm}

%% Colores
\definecolor{azulp}{RGB}{0,90,160}
\definecolor{azulc}{RGB}{220,235,250}
\definecolor{verde}{RGB}{0,130,90}
\definecolor{grisf}{RGB}{248,248,248}
\definecolor{gris}{gray}{0.5}

%% Cajas tcolorbox
\tcbuselibrary{skins,breakable}
\newtcolorbox{cajat}[1][Teoría]{enhanced,breakable,
  colback=azulc,colframe=azulp,
  fonttitle=\bfseries\color{white},coltitle=white,
  attach boxed title to top left={yshift=-2mm,xshift=4mm},
  boxed title style={colback=azulp},
  title={#1},arc=4pt,boxrule=0.8pt,
  left=6pt,right=6pt,top=6pt,bottom=6pt}

\newtcolorbox{cajae}[1][Ejemplo resuelto]{enhanced,breakable,
  colback=grisf,colframe=verde,
  fonttitle=\bfseries\color{white},coltitle=white,
  attach boxed title to top left={yshift=-2mm,xshift=4mm},
  boxed title style={colback=verde},
  title={#1},arc=4pt,boxrule=0.8pt,
  left=6pt,right=6pt,top=6pt,bottom=6pt}

%% Encabezado y pie
\pagestyle{fancy}\fancyhf{}
\lhead{\small\textbf{Departamento de Matemática}}
\rhead{\small 3° Medio --- Matemática}
\lfoot{\small\textit{proscar.cl}}
\rfoot{\small Página \thepage\ de \pageref{LastPage}}
\renewcommand{\headrulewidth}{0.5pt}
\renewcommand{\footrulewidth}{0.4pt}
\setlength{\parindent}{0pt}
\setlength{\parskip}{4pt}

\begin{document}
\begin{center}\large\textbf{Departamento de Matemática}\end{center}
\vspace{-2mm}
\begin{center}
\begin{tabular}{|p{7cm}|p{5cm}|p{3cm}|}
\hline
\textbf{Nombre:} & \textbf{Nivel:} 3° Medio & \textbf{Fecha:} 2026 \\\\
\hline
\textbf{Unidad:} Números Complejos & \multicolumn{2}{l|}{\textbf{Asig.:} Matemática} \\
\hline
\end{tabular}
\end{center}
\vspace{2mm}
\begin{center}\textbf{\Large Guia N°1 --- Números Complejos}\\\normalsize\textit{3° Medio}\\\small OA: \texttt{FG-MATE-3M-OAC-01}\end{center}
\vspace{2mm}
\begin{cajat}[UNIDAD IMAGINARIA Y POTENCIAS DE $i$]
No existe ningún número real $x$ tal que $x^2 = -1$. Definimos la \textbf{unidad imaginaria} $i$:
\[ i^2 = -1 \qquad i = \sqrt{-1} \]

\textbf{Ciclo de potencias (período 4):}
\[ i^1=i,\quad i^2=-1,\quad i^3=-i,\quad i^4=1 \]
En general $i^{4k}=1$ para todo $k\in\mathbb{Z}$.
\end{cajat}

\begin{cajae}[Ejemplo resuelto]
Calcular $i^{235}$.\newline $235 = 4\cdot 58 + 3 \Rightarrow i^{235} = i^3 = -i$
\end{cajae}

\textbf{Ejercicios}

\begin{multicols}{2}
\begin{enumerate}[label=\textbf{\arabic*.},start=1]
\item El valor de $i^{100}$ es:
\begin{tabbing}
\hspace{0.5cm}\=\hspace{3.5cm}\=\kill
\textbf{A)} $1$ \> \textbf{D)} $-i$ \\
\textbf{B)} $-1$ \> \textbf{E)} $0$ \\
\textbf{C)} $i$ \\
\end{tabbing}
\item La expresión $i+i^2+i^3+i^4$ equivale a:
\begin{tabbing}
\hspace{0.5cm}\=\hspace{3.5cm}\=\kill
\textbf{A)} $4i$ \> \textbf{D)} $-1$ \\
\textbf{B)} $1$ \> \textbf{E)} $4$ \\
\textbf{C)} $0$ \\
\end{tabbing}
\item Si $n$ es múltiplo de 4, entonces $i^n$ es:
\begin{tabbing}
\hspace{0.5cm}\=\hspace{3.5cm}\=\kill
\textbf{A)} $i$ \> \textbf{D)} $1$ \\
\textbf{B)} $-i$ \> \textbf{E)} $0$ \\
\textbf{C)} $-1$ \\
\end{tabbing}
\end{enumerate}
\end{multicols}

\vspace{4mm}
\begin{cajat}[NÚMERO COMPLEJO $z = a + bi$]
Un número $z = a+bi$ con $a,b\in\mathbb{R}$ es un \textbf{número complejo}.

\begin{itemize}[leftmargin=1.5em]
  \item $\operatorname{Re}(z)=a$\ (parte real)
  \item $\operatorname{Im}(z)=b$\ (parte imaginaria)
  \item $|z|=\sqrt{a^2+b^2}$\ (módulo)
  \item $\bar{z}=a-bi$\ (conjugado)
\end{itemize}
\end{cajat}

\begin{cajae}[Ejemplo resuelto]
Sea $z=3+4i$. Calcular $|z|$.
\[ |z|=\sqrt{3^2+4^2}=\sqrt{25}=5 \]
\end{cajae}

\textbf{Ejercicios}

\begin{multicols}{2}
\begin{enumerate}[label=\textbf{\arabic*.},start=4]
\item La parte imaginaria de $z = 1-2i$ es:
\begin{tabbing}
\hspace{0.5cm}\=\hspace{3.5cm}\=\kill
\textbf{A)} $-2i$ \> \textbf{D)} $2$ \\
\textbf{B)} $1$ \> \textbf{E)} $-1$ \\
\textbf{C)} $-2$ \\
\end{tabbing}
\item Si $z=5+12i$, entonces $|z|$ es:
\begin{tabbing}
\hspace{0.5cm}\=\hspace{3.5cm}\=\kill
\textbf{A)} $17$ \> \textbf{D)} $169$ \\
\textbf{B)} $13$ \> \textbf{E)} $7$ \\
\textbf{C)} $\sqrt{119}$ \\
\end{tabbing}
\item El conjugado de $z=-3+7i$ es:
\begin{tabbing}
\hspace{0.5cm}\=\hspace{3.5cm}\=\kill
\textbf{A)} $3+7i$ \> \textbf{D)} $-7+3i$ \\
\textbf{B)} $-3-7i$ \> \textbf{E)} $7-3i$ \\
\textbf{C)} $3-7i$ \\
\end{tabbing}
\end{enumerate}
\end{multicols}

\vspace{4mm}
\begin{cajat}[OPERACIONES CON COMPLEJOS]
Sean $z_1=a+bi$ y $z_2=c+di$:

\textbf{Adición:} $z_1+z_2=(a+c)+(b+d)i$

\textbf{Multiplicación:}
\[ z_1\cdot z_2=(ac-bd)+(ad+bc)i \]

\textbf{División} (amplificar por el conjugado):
\[ \frac{z_1}{z_2}=\frac{z_1\cdot\bar{z}_2}{|z_2|^2} \]
\end{cajat}

\begin{cajae}[Ejemplo resuelto]
Calcular $(2+3i)(1-i)$.
\[ (2+3i)(1-i)=2-2i+3i-3i^2=2+i+3=5+i \]
\end{cajae}

\textbf{Ejercicios}

\begin{multicols}{2}
\begin{enumerate}[label=\textbf{\arabic*.},start=7]
\item Si $u=2+3i$ y $v=-5+4i$, entonces $u+v$ es:
\begin{tabbing}
\hspace{0.5cm}\=\hspace{3.5cm}\=\kill
\textbf{A)} $-3+7i$ \> \textbf{D)} $7-3i$ \\
\textbf{B)} $3+7i$ \> \textbf{E)} $2+4i$ \\
\textbf{C)} $-3-7i$ \\
\end{tabbing}
\item El producto $(1+i)(1-i)$ es:
\begin{tabbing}
\hspace{0.5cm}\=\hspace{3.5cm}\=\kill
\textbf{A)} $0$ \> \textbf{D)} $-2$ \\
\textbf{B)} $2i$ \> \textbf{E)} $1+2i$ \\
\textbf{C)} $2$ \\
\end{tabbing}
\item El valor de $\dfrac{1}{i}$ es:
\begin{tabbing}
\hspace{0.5cm}\=\hspace{3.5cm}\=\kill
\textbf{A)} $1$ \> \textbf{D)} $-1$ \\
\textbf{B)} $i$ \> \textbf{E)} $0$ \\
\textbf{C)} $-i$ \\
\end{tabbing}
\end{enumerate}
\end{multicols}

\vspace{4mm}
\vspace{4mm}
\begin{center}\textbf{Soluciones}\end{center}
\begin{center}
\begin{tabular}{|c|c|c|c|c|c|c|c|}
\hline
\textbf{1} & \textbf{2} & \textbf{3} & \textbf{4} & \textbf{5} & \textbf{6} & \textbf{7} & \textbf{8} \\
\hline
A & C & D & C & B & B & A & C \\
\hline
\textbf{9} \\
\hline
C \\
\hline
\end{tabular}
\end{center}

\vfill
\begin{center}\small\color{gris}Generado con \texttt{latex-guia-pdf} $\cdot$ \href{https://proscar.cl}{proscar.cl}\end{center}
\end{document}
